
\textbf{AQUI EMPEZARIA EXPLICANDO COMO ES LA GEOMETRIA Y SPIN STATE DE Pt10 Y Pt5Ge5 ORIGINALES Y COMO CAMBIA EN LOS DERIVADOS A RAIZ DE CAMBIAR LA COMPOSICION}

\textbf{GALDERA: AL CAMBIAR EL METODO HA CAMBIADO LA MULTIPLICIDAD?}

\textbf{EN LA FIG. 2 QUITARIA "REPORTED FRAMEWORK Y FROM Pt10 O FROM Pt5Ge5, PARA QUE QUEDE MAS LIMPIA}

The \ce{Pt10} and \ce{Pt_xGe_{10-x}} clusters considered here
are used to examine how composition, cluster geometry, and
the local Pt/Ge environment affect the affinity of Pt sites for \ce{CO}.
{Figure~\ref{fig:crossed_models} summarizes the two crossed-parent
construction routes used in the primary comparison. We also calculated CO
adsorption starting from the reported global-minimum structures, to examine how
the parent isomer and spin state affect the results. Both sets follow the same
adsorption protocol.}

\begin{figure}[H]
\centering
\begin{tikzpicture}[
  x=1cm,
  y=1cm,
  route/.style={font=\sffamily\bfseries\large, align=center},
  note/.style={font=\sffamily\small, align=center, text=black!70},
  molname/.style={font=\sffamily\bfseries\small, align=center},
  molnote/.style={font=\sffamily\scriptsize, align=center, text=black!70},
  arrow/.style={->, >=stealth, line width=0.7pt, black!85}
]

% Image sizes.  These are intentionally defined here so the layout can be
% adjusted directly from LaTeX.
\def\parentw{0.19\linewidth}
\def\parentwide{0.17\linewidth}
\def\molw{0.15\linewidth}
\def\molwb{0.15\linewidth}

% Molecule images.  Draw these first because the PNGs have white backgrounds.
\node[inner sep=0pt] (pt10parent) at (-3.30,2.75)
  {\includegraphics[width=\parentw]{pt10.png}};
\node[inner sep=0pt] (pt5ge5parent) at (3.45,2.75)
  {\includegraphics[width=\parentwide]{pt5ge5_orig.png}};

\node[inner sep=0pt] (pt6ge4) at (-5.15,-0.10)
  {\includegraphics[width=\molw]{pt6ge4.png}};
\node[inner sep=0pt] (pt5ge5) at (-3.25,-0.10)
  {\includegraphics[width=\molw]{Pt5Ge5_fromPt10Parent.png}};
\node[inner sep=0pt] (pt4ge6) at (-1.15,-0.10)
  {\includegraphics[width=\molw]{pt4ge6.png}};
\node[inner sep=0pt] (pt10frompt5ge5) at (3.45,-0.10)
  {\includegraphics[width=\molwb]{pt10_from_pt5ge5.png}};

% Route arrows.
\draw[arrow] (-3.82,1.18) -- (-4.90,0.80);
\draw[arrow] (-3.30,1.18) -- (-3.25,0.82);
\draw[arrow] (-2.78,1.18) -- (-1.55,0.80);
\draw[arrow] (3.45,1.18) -- (3.45,0.85);

% Headers and labels.  Draw these last so they remain visible over the PNG
% image boxes.
\node[route] at (-3.30,4.58) {\ce{Pt10} parent framework};
\node[note] at (-3.30,4.16) {Ge substitution followed by optimization};
\node[route] at (3.45,4.58) {\ce{Pt5Ge5} parent framework};
\node[note] at (3.45,4.16) {Ge-to-Pt substitution followed by optimization};

\node[molname] at (-3.30,1.62) {\ce{Pt10} parent};
\node[molnote] at (-3.30,1.29) {reported framework};
\node[molname] at (3.45,1.62) {\ce{Pt5Ge5} parent};
\node[molnote] at (3.45,1.29) {reported framework};

\node[molname] at (-5.15,-1.18) {\ce{Pt6Ge4}};
\node[molnote] at (-5.15,-1.51) {from \ce{Pt10}};
\node[molname] at (-3.25,-1.18) {\ce{Pt5Ge5}};
\node[molnote] at (-3.25,-1.51) {from \ce{Pt10}};
\node[molname] at (-1.35,-1.18) {\ce{Pt4Ge6}};
\node[molnote] at (-1.35,-1.51) {from \ce{Pt10}};
\node[molname] at (3.45,-1.18) {\ce{Pt10}};
\node[molnote] at (3.45,-1.51) {from \ce{Pt5Ge5}};

\node[note] at (-0.70,-2.22)
  {Upper row: reported parent frameworks. Lower row: optimized crossed-parent derivatives.};

\end{tikzpicture}

\caption{Crossed-parent strategy used to examine the coupled effects of
composition and structural framework in the ten-atom clusters. Starting from
the reported global-minimum \ce{Pt10} framework, selected Pt atoms were
replaced by Ge and the resulting \ce{Pt6Ge4}, \ce{Pt5Ge5}, and \ce{Pt4Ge6}
structures were optimized. In the reverse construction, the reported
global-minimum \ce{Pt5Ge5} framework was converted into a pure \ce{Pt10}
derivative and then optimized. The snapshots
shown in the upper row correspond to the reported parent frameworks, whereas
the lower row shows the optimized crossed-parent derivatives used in the
present study.}
\label{fig:crossed_models}
\end{figure}

At the PBEPBE/def2-TZVP level, the \ce{Pt10} structure derived from the
\ce{Pt5Ge5} framework is lowest as a triplet, and the quintet is 0.0413 eV
higher in $\Delta E_0$. The selected triplet wavefunction and the seven
distinct triplet \ce{Pt10CO} minima were confirmed to be orbital-stable. The three
Ge-containing derivatives generated from the \ce{Pt10} parent framework favor
the closed-shell singlet over the triplet by 0.176--0.263 eV. Thus, within this
set of models, the Ge-containing clusters have a clearer preference for
the low-spin state than the pure-\ce{Pt10} derivative.

\subsubsection{Bonding patterns in the crossed-parent clusters}

AdNDP was used to compare localized two-center bonding with delocalized
multicenter bonding in the crossed-parent structures.\cite{zubarev:adndp:2008,
tkachenko:adndp2:2019}
For the triplet \ce{Pt10}, the $alpha$ and $beta$ density
matrices were analyzed separately. Each channel gives 30 Pt-centered 1c
elements, with three on every Pt atom. Nineteen 3c elements occur on the same
atomic centers in both channels. Their combined alpha+beta occupation numbers
are 1.984--1.994, so they are assigned as 19 3c--2e elements. The alpha
channel contains two additional 3c elements, % on Pt(2)--Pt(3)--Pt(7) and Pt(5)--Pt(6)--Pt(10), 
each with an occupation number of about 0.994. These are
assigned as two 3c--1e elements and account for the two excess alpha electrons
of the triplet. The large number of high-occupation 3c elements distributed
over the Pt framework supports a delocalized multicenter bonding description.
The complete spin-resolved decomposition and occupation ranges are reported
in Supporting Table~\ref{tab:tfg_adndp_summary}.

This spin-resolved pattern agrees qualitatively with the earlier AdNDP analysis of the
lowest-energy high-spin \ce{Pt10} cluster. That structure was described by 30
one-center pairs, 16 3c--2e bonds, four 3c--1e bonds, and four 4c--1e
bonds.\cite{ugartemendia:2022} The exact numbers are different because the two
\ce{Pt10} structures and spin states are different, but both analyses describe
bonding in pure Pt clusters as mainly delocalized and multicenter.
In contrast, the Ge-containing clusters favor more localized elements. The \ce{Pt5Ge5}
derivative obtained from the \ce{Pt10} framework is dominated by localized
Pt--Ge and Pt--Pt two-center bonds, although one four-center element remains.
The \ce{Pt6Ge4} and \ce{Pt4Ge6} models also favor mainly one- and two-center
elements. The starting geometry affects the detailed bonding pattern, while
Ge substitution favors a more localized and covalent bonding framework.
Representative multicenter and two-center AdNDP elements are shown in
Figure~\ref{fig:tfg_adndp_comparison}. The numerical summaries, occupation
ranges, and search thresholds are given in Supporting Table~\ref{tab:tfg_adndp_summary}.

\begin{figure}[H]
\centering
\begin{tikzpicture}[
  x=1cm,
  y=1cm,
  panel/.style={font=\sffamily\bfseries\small, anchor=north west, fill=white,
    fill opacity=0.86, text opacity=1, inner sep=2pt},
  %label/.style={font=\sffamily\bfseries\small, align=center},
  label/.style={font=\sffamily\small, align=center},
  note/.style={font=\sffamily\scriptsize, align=center, text=black!70,
    text width=0.28\linewidth}
]

% Panel sizes and positions can be tuned directly from here.
\def\adndpw{0.305\linewidth}
\def\yimage{0}
\def\ylabel{-2.42}
\def\ynote{-3.15}

\node[inner sep=0pt] (pt10c) at (-4.35,\yimage)
  {\includegraphics[width=\adndpw]{adndp_pt10_3c.png}};
\node[inner sep=0pt] (pt5ge2c) at (0,\yimage)
  {\includegraphics[width=\adndpw]{adndp_pt5ge_2c2e.png}};
\node[inner sep=0pt] (pt5ge4c) at (4.15,\yimage)
  {\includegraphics[width=\adndpw]{adndp_pt5ge_4c2e.png}};

\node[panel] at ([xshift=-1.72cm,yshift=1.82cm]pt10c.center) {(a)};
\node[panel] at ([xshift=-1.72cm,yshift=1.82cm]pt5ge2c.center) {(b)};
\node[panel] at ([xshift=-1.72cm,yshift=1.82cm]pt5ge4c.center) {(c)};

\node[label] at (-4.15,\ylabel) {\ce{Pt10}};
\node[note] at (-4.15,\ynote) {representative\\3c topology};
\node[label] at (0,\ylabel) {\ce{Pt5Ge5}};
\node[note] at (0,\ynote) {representative Pt--Ge\\2c--2e element};
\node[label] at (4.15,\ylabel) {\ce{Pt5Ge5}};
\node[note] at (4.15,\ynote) {residual\\4c--2e element};

\end{tikzpicture}

\caption{Representative AdNDP elements showing the change from multicenter
metallic bonding in the pure \ce{Pt10} derivative to predominantly localized
two-center bonding in the Ge-containing models: (a) a 3c element of triplet
\ce{Pt10}, (b) a Pt--Ge 2c--2e element of \ce{Pt5Ge5}, and (c) the residual
4c--2e element of \ce{Pt5Ge5}. Panel (a) shows the spatial form of one of the
3c elements identified in the spin-resolved analysis.}
\label{fig:tfg_adndp_comparison}
\end{figure}



\subsubsection[CO adsorption energies]{\ce{CO} adsorption energies}

We next compare \ce{CO} adsorption across the different compositions and parent
structures. Selected spin states are included to check whether they change the
main adsorption patterns. The primary data set is the crossed-parent series described
above. We compare it with a second set obtained by placing \ce{CO} at the
independent Pt sites of the previously reported global-minimum structures of
the same compositions.\cite{ugartemendia:2022} This comparison shows how the
calculated adsorption changes when a different parent geometry or spin state
is used. Figure~\ref{fig:neutral_data_sets} summarizes the relation between
the two sets.


The labels shown in Figure~\ref{fig:tfg_site_numbering} identify the initial
adsorption positions. Because these clusters are flexible, different starting
positions can converge to the same optimized carbonyl minimum after the
cluster and the adsorbate relax.

\begin{figure}[H]
\centering
\begin{tikzpicture}[
  x=1cm,
  y=1cm,
  panel/.style={font=\sffamily\small, anchor=north west, fill=white,
    fill opacity=0.90, text opacity=1, inner sep=2pt},
  title/.style={font=\sffamily\small, align=center},
  note/.style={font=\sffamily\scriptsize, align=center, text=black!70},
  sitelabel/.style={draw=black,
	  circle, fill=white, line width=0.5pt,
	  draw=none,
	  draw opacity=1,
	  fill opacity=0.05,
      text opacity=0.5,
      minimum size=4.8mm, inner sep=0pt, font=\sffamily\bfseries\scriptsize
}
]

% Layout controls.
\def\imgw{0.265\linewidth}
\def\leftx{-3.45}
\def\rightx{3.45}
\def\topy{2.45}
\def\bottomy{-2.75}
\def\titleshift{-2.05cm}

% Bare crossed-parent models.
\node[inner sep=0pt] (pt10cross) at (\leftx,\topy)
  {\includegraphics[width=\imgw]{Pt10CO/pt10_ref.png}};
\node[inner sep=0pt] (pt6ge4cross) at (\rightx,\topy)
  {\includegraphics[width=\imgw]{pt6ge4.png}};
\node[inner sep=0pt] (pt5ge5cross) at (\leftx,\bottomy)
  {\includegraphics[width=\imgw]{Pt5Ge5_fromPt10Parent.png}};
\node[inner sep=0pt] (pt4ge6cross) at (\rightx,\bottomy)
  {\includegraphics[width=\imgw]{pt4ge6.png}};

% Panel markers.
\node[panel] at ([xshift=-1.60cm,yshift=1.52cm]pt10cross.center) {(a)};
\node[panel] at ([xshift=-1.60cm,yshift=1.52cm]pt6ge4cross.center) {(b)};
\node[panel] at ([xshift=-1.60cm,yshift=1.52cm]pt5ge5cross.center) {(c)};
\node[panel] at ([xshift=-1.60cm,yshift=1.52cm]pt4ge6cross.center) {(d)};

% Model labels.
\node[title] at ([yshift=\titleshift]pt10cross.center)
  {\ce{Pt10} from \ce{Pt5Ge5}};
\node[title] at ([yshift=\titleshift]pt6ge4cross.center)
  {\ce{Pt6Ge4} from \ce{Pt10}};
\node[title] at ([yshift=\titleshift]pt5ge5cross.center)
  {\ce{Pt5Ge5} from \ce{Pt10}};
\node[title] at ([yshift=\titleshift]pt4ge6cross.center)
  {\ce{Pt4Ge6} from \ce{Pt10}};

% Pt10 site numbers.
\node[sitelabel] at ([xshift=-.10cm,yshift=-0.02cm]pt10cross.center) {1};
\node[sitelabel] at ([xshift=0.54cm,yshift=0.98cm]pt10cross.center) {3};
\node[sitelabel] at ([xshift=1.15cm,yshift=-0.25cm]pt10cross.center) {4};
\node[sitelabel] at ([xshift=1.18cm,yshift=0.33cm]pt10cross.center) {2};
\node[sitelabel] at ([xshift=-0.39cm,yshift=1.26cm]pt10cross.center) {10};
\node[sitelabel] at ([xshift=-1.30cm,yshift=0.66cm]pt10cross.center) {8};
\node[sitelabel] at ([xshift=-1.21cm,yshift=-0.33cm]pt10cross.center) {7};
\node[sitelabel] at ([xshift=-.53cm,yshift=-1.27cm]pt10cross.center) {6};
\node[sitelabel] at ([xshift=.66cm,yshift=-0.90cm]pt10cross.center) {5};
\node[sitelabel] at ([xshift=.12cm,yshift=.57cm]pt10cross.center) {9};

% Pt6Ge4 representative adsorption-site classes. Positions are tied to the
% selected Jmol view and can be fine-tuned independently.
\node[sitelabel] at ([xshift=-0.83cm,yshift=0.30cm]pt6ge4cross.center) {1};
\node[sitelabel] at ([xshift=-0.03cm,yshift=0.27cm]pt6ge4cross.center) {2};
\node[sitelabel] at ([xshift=0.82cm,yshift=0.32cm]pt6ge4cross.center) {4};
\node[sitelabel] at ([xshift=-0.02cm,yshift=-1.05cm]pt6ge4cross.center) {9};
\node[sitelabel] at ([xshift=-0.92cm,yshift=-1.1cm]pt6ge4cross.center) {7};
\node[sitelabel] at ([xshift=1.02cm,yshift=-1.05cm]pt6ge4cross.center) {5};

% Pt5Ge5 site-class numbers.
\node[sitelabel] at ([xshift=-.80cm,yshift=.30cm]pt5ge5cross.center) {6};
\node[sitelabel] at ([xshift=.68cm,yshift=.30cm]pt5ge5cross.center) {8};
\node[sitelabel] at ([xshift=-0.1cm,yshift=0.0cm]pt5ge5cross.center) {4};
\node[sitelabel] at ([xshift=-0.96cm,yshift=-0.85cm]pt5ge5cross.center) {7};
\node[sitelabel] at ([xshift=0.82cm,yshift=-0.85cm]pt5ge5cross.center) {10};

% Pt4Ge6 adsorption sites.
\node[sitelabel] at ([xshift=-0.80cm,yshift=0.39cm]pt4ge6cross.center) {1};
\node[sitelabel] at ([xshift=0.68cm,yshift=0.39cm]pt4ge6cross.center) {2};
\node[sitelabel] at ([xshift=-0.42cm,yshift=-0.46cm]pt4ge6cross.center) {5};
\node[sitelabel] at ([xshift=0.17cm,yshift=-1.08cm]pt4ge6cross.center) {7};

\end{tikzpicture}


\caption{Numbering of the Pt atoms used to place \ce{CO} and calculate the
\ce{CO} binding energies. Panel (a) shows the \ce{Pt10} model derived
from the \ce{Pt5Ge5} parent framework. Panels (b)--(d) show the \ce{Pt6Ge4},
\ce{Pt5Ge5}, and \ce{Pt4Ge6} derivatives,
respectively, generated from the \ce{Pt10} parent framework. The labels retain
the atom indices of the corresponding parent structure and identify the
starting adsorption sites.}
\label{fig:tfg_site_numbering}
\end{figure}


The ten triplet starting-site optimizations produce seven unique
\ce{Pt10CO} minima, which are collected in
Figure~\ref{fig:tfg_pt10co_series}. The paired labels 1/9, 3/5, and 8/10
indicate starting structures that converge to the same relaxed minimum. These
calculations provide a systematic set of trajectories from the selected
\ce{Pt10} isomer.

\begin{figure}[H]
\centering
\begin{tikzpicture}[
  x=1cm,
  y=1cm,
  structure/.style={inner sep=0pt},
  sitelabel/.style={font=\sffamily\scriptsize, align=center},
  datalabel/.style={font=\sffamily\tiny, align=center, text=black!75}
]

% Global layout controls.
\def\molw{0.22\linewidth}
\def\rowone{1.75}
\def\rowtwo{-1.75}
\def\xone{-4.35}
\def\xtwo{-1.45}
\def\xthree{1.45}
\def\xfour{4.35}
\def\labelshift{-1.48cm}
\def\datashift{-1.77cm}

% Seven unique minima obtained from the ten original TFG starting sites.
\node[structure] (site1) at (\xone,\rowone)
  {\includegraphics[width=\molw]{Pt10CO/Pt1_from_1xyzUnai.png}};
\node[structure] (site2) at (\xtwo,\rowone)
  {\includegraphics[width=\molw]{Pt10CO/Pt2_from_2xyzUnai.png}};
\node[structure] (site5) at (\xthree,\rowone)
  {\includegraphics[width=\molw]{Pt10CO/Pt5_from_5xyzUnai.png}};
\node[structure] (site4) at (\xfour,\rowone)
  {\includegraphics[width=\molw]{Pt10CO/pt4_from_4xyzUnai.png}};

\node[structure] (site6) at (-4.0,\rowtwo)
  {\includegraphics[width=\molw]{Pt10CO/Pt6_from_6xyzUnai.png}};
\node[structure] (site7) at (0,\rowtwo)
  {\includegraphics[width=\molw]{Pt10CO/Pt7_from_7xyzUnai.png}};
\node[structure] (site8) at (4.0,\rowtwo)
  {\includegraphics[width=\molw]{Pt10CO/Pt8_from_8xyzUnai.png}};

% Starting-site labels and binding energies of the final minima.
\node[sitelabel] at ([yshift=\labelshift]site1.center) {Sites 1/9};
\node[datalabel] at ([yshift=\datashift]site1.center) {$-2.156$ eV};
\node[sitelabel] at ([yshift=\labelshift]site2.center) {Site 2};
\node[datalabel] at ([yshift=\datashift]site2.center) {$-2.493$ eV};
\node[sitelabel] at ([yshift=\labelshift]site5.center) {Sites 3/5};
\node[datalabel] at ([yshift=\datashift]site5.center) {$-2.590$ eV};
\node[sitelabel] at ([yshift=\labelshift]site4.center) {Site 4};
\node[datalabel] at ([yshift=\datashift]site4.center) {$-2.579$ eV};

\node[sitelabel] at ([yshift=\labelshift]site6.center) {Site 6};
\node[datalabel] at ([yshift=\datashift]site6.center) {$-2.657$ eV};
\node[sitelabel] at ([yshift=\labelshift]site7.center) {Site 7};
\node[datalabel] at ([yshift=\datashift]site7.center) {$-2.493$ eV};
\node[sitelabel] at ([yshift=\labelshift]site8.center) {Sites 8/10};
\node[datalabel] at ([yshift=\datashift]site8.center) {$-2.650$ eV};

\end{tikzpicture}

\caption{Seven unique optimized triplet \ce{Pt10CO} minima obtained from the
ten initial Pt adsorption sites of the \ce{Pt10} derivative. Site
labels follow Figure~\ref{fig:tfg_site_numbering}; paired labels identify
initial structures that converge to the same final minimum. The corresponding
triplet $BE(\ce{CO})$ values are given below each structure. {The
analogous quintet calculations produced five distinct minima; their relative
and binding energies are also reported together with the triplet results in
Table~\ref{tab:tfg_adsorption_descriptors}.}}
\label{fig:tfg_pt10co_series}
\end{figure}

The optimized carbonyl structures of the three Ge-containing compositions are
shown in Figure~\ref{fig:tfg_ptgeco_series}.

\begin{figure}[H]
\centering
\begin{tikzpicture}[
  x=1cm,
  y=1cm,
  structure/.style={inner sep=0pt},
  rowlabel/.style={font=\sffamily\small\bfseries, align=center},
  sitelabel/.style={font=\sffamily\scriptsize, align=center},
  datalabel/.style={font=\sffamily\tiny, align=center, text=black!75}
]

% Global layout controls.
\def\molwA{0.17\linewidth}
\def\molwB{0.13\linewidth}
\def\molwC{0.15\linewidth}
\def\molw{0.215\linewidth}
\def\rowone{3.55}
\def\rowtwo{0.00}
\def\rowthree{-3.55}
\def\xone{-4.65}
\def\xtwo{-1.55}
\def\xthree{1.55}
\def\xfour{4.65}
\def\xmidleft{-3.10}
\def\xmid{0.00}
\def\xmidright{3.10}
\def\labelshift{-1.47cm}
\def\datashift{-1.73cm}

% Pt6Ge4CO adsorption sites.
\node[structure] (p6g4s1) at (\xone,\rowone)
  {\includegraphics[width=\molwB]{pt6ge4_site01.png}};
\node[structure] (p6g4s2) at (\xtwo,\rowone)
  {\includegraphics[width=\molwB]{pt6ge4_site02.png}};
\node[structure] (p6g4s7) at (\xthree,\rowone)
  {\includegraphics[width=\molwA]{pt6ge4_site07.png}};
\node[structure] (p6g4s9) at (\xfour,\rowone)
  {\includegraphics[width=\molwA]{pt6ge4_site09.png}};

\node[rowlabel] at (0,{\rowone+1.48}) {\ce{Pt6Ge4CO}};
\node[sitelabel] at ([yshift=\labelshift]p6g4s1.center) {Site 1/4};
\node[datalabel] at ([yshift=\datashift]p6g4s1.center) {$-2.235$ eV};
\node[sitelabel] at ([yshift=\labelshift]p6g4s2.center) {Site 2};
\node[datalabel] at ([yshift=\datashift]p6g4s2.center) {$-2.246$ eV};
\node[sitelabel] at ([yshift=\labelshift]p6g4s7.center) {Site 5/7};
\node[datalabel] at ([yshift=\datashift]p6g4s7.center) {$-1.567$ eV};
\node[sitelabel] at ([yshift=\labelshift]p6g4s9.center) {Site 9};
\node[datalabel] at ([yshift=\datashift]p6g4s9.center) {$-2.245$ eV};

% Pt5Ge5CO adsorption-site classes. Parent-Pt indices are retained; paired
% labels denote equivalent sites represented by the same optimized structure.
\node[structure] (p5g5s1) at (\xmidleft,\rowtwo)
  {\includegraphics[width=\molw]{Pt5Ge5CO/pt5ge5_1_from_1xyz.png}};
\node[structure] (p5g5s2) at (\xmid,\rowtwo)
  {\includegraphics[width=\molw]{Pt5Ge5CO/Pt5Ge5_2_from_2xyzUnai.png}};
\node[structure] (p5g5s3) at (\xmidright,\rowtwo)
  {\includegraphics[width=\molw]{Pt5Ge5CO/Pt5Ge5_3_from_3xyzUnai.png}};

\node[rowlabel] at (0,{\rowtwo+1.48}) {\ce{Pt5Ge5CO}};
\node[sitelabel] at ([yshift=\labelshift]p5g5s1.center) {Sites 6/8};
\node[datalabel] at ([yshift=\datashift]p5g5s1.center) {$-2.076$ eV};
\node[sitelabel] at ([yshift=\labelshift]p5g5s2.center) {Sites 7/10};
\node[datalabel] at ([yshift=\datashift]p5g5s2.center) {$-1.661$ eV};
\node[sitelabel] at ([yshift=\labelshift]p5g5s3.center) {Site 4};
\node[datalabel] at ([yshift=\datashift]p5g5s3.center) {$-1.342$ eV};

% Pt4Ge6CO adsorption sites.
\node[structure] (p4g6s1) at (\xone,\rowthree)
  {\includegraphics[width=\molwC]{pt4ge6_site01.png}};
\node[structure] (p4g6s2) at (\xtwo,\rowthree)
  {\includegraphics[width=\molwB]{pt4ge6_site02.png}};
\node[structure] (p4g6s5) at (\xthree,\rowthree)
  {\includegraphics[width=\molwB]{pt4ge6_site05.png}};
\node[structure] (p4g6s7) at (\xfour,\rowthree)
  {\includegraphics[width=\molwC]{pt4ge6_site07.png}};

\node[rowlabel] at (0,{\rowthree+1.48}) {\ce{Pt4Ge6CO}};
\node[sitelabel] at ([yshift=\labelshift]p4g6s1.center) {Site 1};
\node[datalabel] at ([yshift=\datashift]p4g6s1.center) {$-1.052$ eV};
\node[sitelabel] at ([yshift=\labelshift]p4g6s2.center) {Site 2};
\node[datalabel] at ([yshift=\datashift]p4g6s2.center) {$-1.305$ eV};
\node[sitelabel] at ([yshift=\labelshift]p4g6s5.center) {Site 5};
\node[datalabel] at ([yshift=\datashift]p4g6s5.center) {$-1.459$ eV};
\node[sitelabel] at ([yshift=\labelshift]p4g6s7.center) {Site 7};
\node[datalabel] at ([yshift=\datashift]p4g6s7.center) {$-1.458$ eV};

\end{tikzpicture}

\caption{Optimized carbonyl structures of the \ce{Pt6Ge4}, \ce{Pt5Ge5},
and \ce{Pt4Ge6} derivatives generated from the \ce{Pt10} parent framework
(top, middle, and bottom rows, respectively). Site labels retain the Pt atom
numbering of the parent structure shown in Figure~\ref{fig:tfg_site_numbering}.
For \ce{Pt6Ge4CO}, the four structures represent sites or site classes
$2$, $1/4$, $5/7$, and $9$ identified in the parent geometry. The paired
labels denote equivalent Pt positions represented by one optimized structure.
For \ce{Pt5Ge5CO}, three representative Pt positions are shown. The paired
positions 6/8 and 7/10 were considered equivalent and are represented by one
optimized structure each. For \ce{Pt4Ge6CO}, sites 5 and 7 converge to
effectively the same carbonyl minimum and are shown once as 5/7. The
corresponding $BE(\ce{CO})$ values are given
below each structure.}
\label{fig:tfg_ptgeco_series}
\end{figure}
